\section{The Events of 1917 in Their Nearest Witnesses}

\subsection{The seers and the sources}

The three children were Lucia de Jesus dos Santos (born in Aljustrel,
parish of F\'atima, 28 March 1907, registered as 22 March), and her
first cousins Francisco Marto (born 1908) and Jacinta Marto (born
1910), children of shepherd families of Aljustrel.\endnote{Birth and
family data: the editorial notes to Doc.~1 of the \work{Sele\c{c}\~ao
de documentos} (the parish priest's first interrogation), which record
Lucia's baptismal register date and family; the Sanctuary's
``Chronology of the Three Seers'' page,
\url{https://www.fatima.pt/en/pages/chronology-of-the-three-seers},
which explains the 22/28 March discrepancy from Lucia's Fifth Memoir;
and the Dicastery for the Causes of Saints' decree on Lucia's heroic
virtues, 22 June 2023 (References). The ages printed under the
photograph in \work{O S\'eculo} (10, 9, 7) are the newspaper's.} In
1917 they reported six apparitions of a Lady at the Cova da Iria,
about 2.5~km from F\'atima, on 13 May, 13 June, 13 July, 19 August
(after their detention prevented the expected 13 August encounter),
13 September, and 13 October.

The evidentiary strata must be kept distinct, because they were
written at very different distances from the events. Nearest are the
contemporaneous records of 1917: the parish priest of F\'atima, Manuel
Marques Ferreira, began interrogating the children within about two
weeks of the first apparition and after each subsequent one; the
Lisbon press reported from July onward; and Canon Manuel Nunes
Formig\~ao interrogated the children repeatedly from September. These
records were critically edited by the Sanctuary of F\'atima in the
fifteen-volume \work{Documenta\c{c}\~ao Cr\'itica de F\'atima}
(1992--2013); this study uses the Sanctuary's one-volume critical
selection of the 1917--1930 documents, which it distributes freely
(References). Farthest are Lucia's memoirs of 1935--1941 and later
writings, treated in \S4. The contemporaneous records are the
strongest evidence for what was publicly said, asked, and observed in
1917; they do not by themselves prove a supernatural origin, and
nothing below asserts one.

\subsection{May to September in the 1917 record}

The parish priest's first interrogation of Lucia (about 27 May 1917)
already fixes the essential shape of the reported first encounter: a
Lady seen above a holm oak, a request to return, a promise of
heaven.\endnote{\work{Sele\c{c}\~ao}, Doc.~1 (1917-05-c.27),
publishing DCF~I, Doc.~1. The interrogation records of the parish
priest are summaries made by him, not stenographic transcripts.} After
the July apparition he interrogated Lucia and Jacinta on 14 July: the
Lady asked that they return each 13th and that people ``pray the
ter\c{c}o [five-decade Rosary] to Our Lady of the Rosary, that she may
abate the war, for only she can help it''; asked by Lucia for a
miracle so that all would believe, the Lady answered, in the record's
words, \term{``Daqui a tr\^es meses farei ent\~ao com que todos
acreditem''}---within three months I shall make all
believe.\endnote{\work{Sele\c{c}\~ao}, Doc.~3 (1917-07-14),
publishing DCF~I, Doc.~3. The promise of a sign three months ahead is
therefore attested in writing three months \emph{before} 13 October
---the single most important contemporaneous control on the October
crowd.} The same summer records attest that the children claimed a
\emph{secret} they would not reveal: the questioning after 13 August
turns on it, and Formig\~ao's autumn interrogations press all three
children about it without obtaining its content.\endnote{Secret
attested in 1917: e.g.\ \work{Sele\c{c}\~ao}, Doc.~120's account of
the 13 August detention (the administrator sought ``to tear from them
the secret''), and Formig\~ao's September--October interrogations
(Docs.~8--12, 18), which ask each child about the secret and record
their refusal. The 1917 record proves that a secret was claimed; it
does not contain the content later written in 1941 and 1944.}

On 13 August the children never reached the Cova: the administrator of
the county of Vila Nova de Our\'em, Artur de Oliveira Santos, took
them by carriage to Our\'em, held them, interrogated them under
threats, and released them on the 15th; a crowd at the Cova that day
reported unusual atmospheric effects but no vision. The children
reported a compensating apparition at Valinhos, near Aljustrel, on
19 August---a date fixed in Lucia's later Fourth Memoir; the
apparition itself is attested in the parish priest's interrogation of
27 August 1917.\endnote{Detention: \work{Sele\c{c}\~ao}, Doc.~120 (the
diocesan commission's 1930 report, drawing on the sworn 1924
interrogation) and the 1917 press; Valinhos: the Sanctuary's
``Narrative of the Apparitions'' page, which prints the Fourth Memoir
account and notes that the bracketed portion comes from the parish
priest's interrogation of 27 August 1917 (DCF~I, p.~17). The
\emph{day} 19 August rests on the memoir stratum; the fact of an
August apparition away from the Cova is contemporaneous.} September
13 followed the established pattern before a growing crowd, with the
Lady---in the reports---promising that in October ``Our Lord, Our
Lady of Sorrows and Our Lady of Carmel, and Saint Joseph with the
Child Jesus'' would come, and that a miracle would be worked.

\subsection{13 October and the solar phenomenon}

The promised sign drew an enormous crowd to a bare upland in heavy
rain. The most quoted description was written by a witness with no
stake in the Church's favor: Avelino de Almeida (1873--1932), special
correspondent---\term{enviado especial}---of \work{O S\'eculo}, the
Lisbon daily the Sanctuary's critical edition identifies as the
``daily of the Portuguese Republican Party,'' organ of the political
world then closing churches and jailing priests; Almeida's own
pre-report of 13 October had treated the affair with open irony. His
signed report, published 15 October under the title ``Como o sol
bailou ao meio dia em F\'atima'' (How the sun danced at midday at
F\'atima), estimates the crowd---citing ``dispassionate calculations
of cultured persons wholly foreign to mystical influences''---at
thirty to forty thousand, and describes the moment this way: the rain
stopped, the sun appeared ``like a disc of dull silver'' that could be
looked at without strain, and then, ``before the astonished eyes of
that people\ldots\ the sun trembled, the sun made brusque movements,
never seen, outside all cosmic laws---the sun `danced,' in the typical
expression of the peasants.''\endnote{\work{Sele\c{c}\~ao}, Doc.~13
(1917-10-15), publishing DCF~III-1, Doc.~87, at the volume's printed
pages 72--78; the crowd estimate and the quoted sentences are at
pp.~75--76 (editorial translation; Portuguese: ``o sol tremeu, o sol
teve nunca vistos movimentos bruscos fora de todas as leis c\'osmicas
--- o sol `bailou', segundo a t\'ipica express\~ao dos camponeses'').
Almeida died in 1932; the underlying 1917 newspaper text is in the
public domain.} He closed by handing the question on, not by settling
it: ``It remains for the competent to pronounce on the macabre dance
of the sun which today, at F\'atima, made hosannas burst from the
breasts of the faithful and naturally impressed---so trustworthy
persons assured me---freethinkers and other people without religious
concerns.'' Two weeks later, in the same paper's weekly
\work{Ilustra\c{c}\~ao Portuguesa} (29 October), answering a friend
who had lost his faith and wanted ``an unsuspect testimony,'' Almeida
held his ground: ``Miracle, as the people shouted; natural
phenomenon, as the learned say? I do not now care to know, but only
to tell you what I saw\ldots\ The rest is for Science and the
Church.''\endnote{\work{Sele\c{c}\~ao}, Doc.~21 (1917-10-29),
publishing DCF~III-1, Doc.~228, printed pages 114--116 (editorial
translation).}

The honest boundaries of this dossier are as important as its
strength. The phenomenon was reported by believers and unbelievers,
near the Cova and at some distance; descriptions vary (rotation,
color changes, an apparent fall toward the earth); some present
reported seeing nothing unusual; no astronomical observatory recorded
any solar anomaly, so whatever happened was local to the observers.
The 1930 pastoral letter itself argued the phenomenon was witnessed
by tens of thousands ``of all categories and social classes'' and
was ``not natural,'' but the operative clause it issued judges the
children's \emph{visions} worthy of belief; the Church has never
defined a physical mechanism, and this study does not certify one,
adjudicate every witness, or resolve the natural, psychological, and
miraculous hypotheses that have been proposed.\endnote{The pastoral
letter's argumentation about the solar phenomenon precedes its
operative clauses (\work{Sele\c{c}\~ao}, Doc.~133). The negative
astronomical control (no registered global anomaly) is common ground
in the literature and follows from the event's local visibility; it
is stated here as a bound, not as a refutation.}

What the children themselves reported seeing on 13 October---the Lady
as Our Lady of the Rosary asking for a chapel and for the daily
Rosary, then tableaux beside the sun of Saint Joseph with the Child,
Our Lord, Our Lady of Sorrows, and Our Lady of Carmel---is recorded
in the parish priest's interrogation of 16 October and, at greater
length, in the memoirs.\endnote{\work{Sele\c{c}\~ao}, Doc.~14
(1917-10-16), publishing DCF~I, Doc.~6; Fourth Memoir account as
printed by the Sanctuary's ``Narrative of the Apparitions'' page.}

\subsection{A preserved discrepancy: the October war statement}

The same 16 October interrogation records Lucia reporting the Lady's
words as: \term{``a guerra acaba ainda hoje; esperem c\'a pelos seus
militares, muito breve''}---the war ends this very day; expect your
soldiers here very soon. The war did not end that day; it ended in
November 1918. Formig\~ao pressed exactly this point on 19 October:
``But look, the war is still going on!\ldots\ How is this explained,
if Our Lady said the war ended that day?'' Lucia: ``I do not know. I
only know that I heard her say that the war would end on the 13th. I
do not know anything else.'' And, pressed with her own earlier
statement that the end would come ``soon after'': ``I no longer
remember well how she said it. She may have said that; I do not know.
Perhaps I did not understand the Lady well.''\endnote{Interrogation
of 16 October: \work{Sele\c{c}\~ao}, Doc.~14; Formig\~ao's
interrogation (redacted after 19 October): Doc.~18, publishing DCF~I,
Doc.~16, questions 1--6 (editorial translation). The later Fourth
Memoir stratum gives the October words as ``the war is going to
end''; the 1941 text and the 1917 records are not verbally
identical.}

The discrepancy is published here because it is consequential and
because both honest options remain open within the Church's own
frame: either the ten-year-old, questioned in a chaos of light and
crowd, misreported or misremembered a sentence, or the sentence was
conditioned or temporally compressed in ways the hearer could not
carry. The 2000 theological commentary supplies the controlling
category: interior visions are received ``in the modes of
representation and consciousness available'' to the seer, and are
``never simple `photographs' of the other world'' (\S5). A
worthy-of-belief judgment is not a warranty of verbal inerrancy, and
the 1930 act nowhere claims one. What the discrepancy does not
license is silently correcting the 1917 record to match the memoirs,
or treating the whole dossier as thereby discredited: the bishop
judged in 1930 with these same documents before his commission.
